% ============================================================
% 分片：s07-references  对应原文 PDF 第 87--92 页（印刷页码 87--92）
% 负责章节：REFERENCES（文献编号 1--53）
% 译者：t07-s34   审校状态：已审校
% 审校：t01-intro 已完成逐条核对（文献 1--53 逐条对照 pages/page-87..92.png；修正 7 处批注/引用，书目文字零改动）
% 边界说明（Lead 裁定）：p87 顶部（§6.2 收尾句 + 正文一段 + Acknowledgements 致谢段）归 s06b；
%   本分片自 p87 的 “REFERENCES” 标题开始，至 p92 末条目 53 结束。
% 体例：书目信息逐字保留（仅修正 TeX 连字（fi、ff、ffi 等）及跨行断字连字符，并按标准英文引号
%   排配原文中的引号）；标题一律不译；书名/丛书名保持原文斜体，期刊卷号保持原文粗体。
%   “REFERENCES” 标题由 thebibliography 自动生成（main.tex 已设 \refname 为“参考文献”），故不另加 \section*。
% ============================================================

% ==== 原文 PDF 第 87 页（印刷页 87）====
\begin{thebibliography}{99}

% [Tutor Note: 文献 [1--3] 是 C-Y/L-G 对应的原始文献：Martinec 把 Landau-Ginzburg 模型读作紧化中的“灾变”（catastrophe），Vafa--Warner 用灾变理论分类 N=2 共形场论，Greene--Vafa--Warner 则首次系统论证 Calabi-Yau 与 Landau-Ginzburg 属同一 N=2 理论的不同相；本文 §3 的“相”图像即建立在这三篇之上。]
% [原文疑误: 第 1 条作者 “E.Martinec” 原文缩写点号后无空格；按“书目逐字保留”照录，未补空格。]
% [文献 1: Martinec，Landau–Ginzburg 作为紧化“灾变”的早期文献，C-Y/L-G 对应的源头之一。]
\bibitem{ref1} E.Martinec, ``Criticality, Catastrophes, and Compactifications,'' in \textit{Physics And Mathematics Of Strings}, ed L. Brink, D. Friedan, and A. M. Polyakov (World Scientific, 1990).

\bibitem{ref2} C. Vafa and N. Warner, ``Catastrophes And The Classification Of Conformal Field Theories,'' Phys. Lett. \textbf{218B} (1989) 51.

% [文献 3: Greene–Vafa–Warner，Calabi–Yau 与 Landau–Ginzburg 对应的奠基论文，本文主题的直接来源。]
\bibitem{ref3} B. Greene, C. Vafa, and N. Warner, ``Calabi-Yau Manifolds And Renormalization Group Flows,'' Nucl. Phys. \textbf{B324} (1989) 371.

% [文献 4: 手征环（chiral ring）概念的奠基文献；§3.3 讨论的 ac/cc 手征环即出自此。]
\bibitem{ref4} W. Lerche, C. Vafa, and N. Warner, ``Chiral Rings In $N=2$ Superconformal Theory,'' Nucl. Phys. \textbf{B324} (1989) 427.

\bibitem{ref5} S. Cecotti, L. Girardello, and A. Pasquinucci, ``Nonperturbative Aspects And Exact Results In The $N=2$ Landau-Ginzburg Models,'' Nucl. Phys. \textbf{B328} (1989) 701, ``Singularity Theory And $N=2$ Supersymmetry,'' Int. J. Mod. Phys. \textbf{A6} (1991) 2427.

\bibitem{ref6} S. Cecotti, ``$N=2$ Landau-Ginzburg vs. Calabi-Yau $\sigma$-Models: Non-Perturbative Aspects,'' Int. J. Mod. Phys. \textbf{A6} (1991) 1749.

% [Tutor Note: 文献 [7]（AGM）讨论“多重镜像流形”（multiple mirror manifolds）与弦论中的拓扑改变：他们找到拓扑不同但双有理等价、且共享同一镜像的 Calabi--Yau 流形对；§5.5 说明这类局部拓扑跃迁正是 Mori 理论意义下的 flop（[44] 综述了相应的双有理几何），故 [7] 为“物理过程可改变时空拓扑”提供了首批具体证据。]
% [文献 7: Aspinwall–Greene–Morrison，镜像流形与拓扑改变，§5.5 的物理来源。]
\bibitem{ref7} P. Aspinwall, B. Greene, and D. Morrison, ``Multiple Mirror Manifolds And Topology Change In String Theory,'' IAS preprint (1993).

% ==== 原文 PDF 第 88 页（印刷页 88）====
\label{op:88}
\bibitem{ref8} T. Hubsch, ``Chameleonic $\sigma$-Models,'' Phys. Lett. \textbf{247B} (1990) 317, ``How Singular A Space Can Superstrings Thread?'' Mod. Phys. Lett. \textbf{A6} (1991) 207; J. Gates and T. Hubsch, ``Unidexterous Locally Supersymmetric Actions For Calabi-Yau Compactifications,'' Phys. Lett. \textbf{226B} (1989) 100, ``Calabi-Yau Heterotic Strings And Supersymmetric Sigma Models,'' Nucl. Phys. \textbf{B343} (1990) 741.

% [Tutor Note: 文献 [9]（Thaddeus）用稳定对（stable pairs）与线性系给出 $SU(2)$ WZW 模型 Verlinde 公式的代数几何证明；本文 §3.4 处理涡旋方程 (3.47) 时所用的“复商 = 辛商”思想（Kempf--Ness 定理，§4 中作初等说明）与 [9] 的稳定对语言一脉相承，是本文把 C-Y/L-G 相结构翻译成 GIT 语言的关键数学参照。]
% [文献 9: Thaddeus 的稳定对与 Verlinde 公式；本文 §3.4 的涡旋方程与 §4 的商构造（GIT）思想与此相通。]
\bibitem{ref9} M. Thaddeus, ``Stable Pairs, Linear Systems, And The Verlinde Formula,'' MSRI preprint (1992).

\bibitem{ref10} J. Wess and J. Bagger, \textit{Supersymmetry And Supergravity}, Princeton University Press (second edition, 1992).

% [文献 11: Gates–Grisaru–Rocek–Siegel 的超空间标准参考书，(2,2) 超场计算的基本工具。]
\bibitem{ref11} S. J. Gates, Jr., M. T. Grisaru, M. Rocek, and W. Siegel, \textit{Superspace, or One Thousand And One Lessons In Supersymmetry}, (Benjamin-Cummings, 1983).

\bibitem{ref12} P. G. O. Freund, \textit{Introduction To Supersymmetry} (Camb. Univ. Press, 1986).

\bibitem{ref13} P. West, \textit{Introduction To Supersymmetry And Supergravity} (World Scientific, 1986).

\bibitem{ref14} S. P. Misra, \textit{Introduction To Supersymmetry And Supergravity} (Wiley Eastern Limited, 1992).

\bibitem{ref15} J. Gates, ``Superspace Formulation Of New Non-Linear Sigma Models,'' Nucl. Phys. \textbf{B238} (1984) 349.

% [文献 16: Gates–Hull–Rocek 的 twisted multiplets，即“扭手征多重态”的原始文献之一；本文 (2,2) 理论中 $\sigma$ 的扭手征性质与此相关。]
\bibitem{ref16} S. J. Gates, C. M. Hull, and M. Rocek, ``Twisted Multiplets And New Supersymmetric Nonlinear $\sigma$-Models,'' Nucl. Phys. \textbf{B248} (1984) 157

\bibitem{ref17} T. Buscher, U. Lindstrom, and M. Rocek, ``New Supersymmetric $\sigma$-Models With Wess-Zumino Term,'' Phys. Lett. \textbf{202B} (1988) 94.

\bibitem{ref18} N. J. Hitchin, A. Karlhede, U. Lindstrom, and M. Rocek, ``Hyper-Kahler Metrics And Supersymmetry,'' Commun. Math. Phys. \textbf{108} (1987) 535.

% ==== 原文 PDF 第 89 页（印刷页 89）====
\label{op:89}
\bibitem{ref19} M. Rocek, ``Modified Calabi-Yau Manifolds With Torsion,'' in S.-T. Yau, ed., \textit{Essays On Mirror Manifolds} (International Press, 1992).

\bibitem{ref20} M. Rocek and E. Verlinde, ``Duality, Quotients, and Currents,'' Nucl. Phys. \textbf{B373} (1992) 630.

% [文献 21: $\mathbf{CP}^{n-1}$ 模型的大 $N$ 有效作用量；§2 讨论扭超势（twisted superpotential）时与 [20] 一起被引作密切相关的工作。]
\bibitem{ref21} A. D'Adda, A. C. Davis, P. DiVecchia, and P. Salomonson, ``An Effective Action For The $\mathbf{CP}^{n-1}$ Model,'' Nucl. Phys. \textbf{B222} (1983) 45.

\bibitem{ref22} T. Hubsch, ``Of Marginal Kinetic Terms And Anomalies,'' Mod. Phys. Lett. \textbf{A6} (1991) 1553

\bibitem{ref23} S. Coleman, ``More On The Massive Schwinger Model,'' Ann. Phys. (NY) \textbf{101} (1976) 239.

\bibitem{ref24} E. Witten, ``Mirror Manifolds And Topological Field Theory,'' in S.-T. Yau, ed., \textit{Essays On Mirror Manifolds} (International Press, 1992).

\bibitem{ref25} C. Vafa and S. Cecotti, ``Exact Results For Supersymmetric Sigma Models,'' Phys. Rev. Lett. \textbf{68} (1992) 903.

\bibitem{ref26} A. Schellekens and N. Warner, ``Anomaly Cancellation And Self-Dual Lattices,'' Phys. Lett. \textbf{181B} (1986) 339; K. Pilch, A. Schellekens and N. Warner, ``Anomalies, Characters, and Strings,'' Nucl. Phys. \textbf{B287} (1987) 317.

% [文献 27: Witten 关于椭圆亏格与拓扑场论的经典文献；§3.3 的椭圆亏格 (3.29) 与其指标解释均源于此。]
\bibitem{ref27} E. Witten, ``Elliptic Genera And Quantum Field Theory,'' Commun. Math. Phys. \textbf{109} (1987) 525, ``The Index Of The Dirac Operator In Loop Space,'' in P. Landwebber, cited below.

% [存疑: 第 27、28 条 “Landwebber” 通常拼作 Landweber（Peter Landweber）；第 32 条 “Daskapoulos” 通常拼作 Daskalopoulos；第 38 条 “M. Aubin” 通行作者名为 Mich\`ele Audin。以上均照录原文，未改。另：原文 [16]、[22] 末尾缺句点，[17] 收尾误排为单引号，[26] 缺收尾双引号，[27] 第二篇标题起始引号误排为 ‘’；本稿按译者声明的引号体例统一为标准英文双引号，句读照录原文位置，书目文字未作改动。又：原文 [47] 的作者依次为 Candelas, Green, Parke, de la Ossa，与该文通行署名顺序（Candelas, de la Ossa, Green, Parke）不同，此处照录原文。]
\bibitem{ref28} P. Landwebber, ed., \textit{Elliptic Curves And Modular Forms In Algebraic Topology} (Springer-Verlag, 1988).

\bibitem{ref29} E. Witten, ``Topological Sigma Models,'' Commun. Math. Phys. \textbf{118} (1988) 411.

\bibitem{ref30} T. Eguchi and S.-K. Yang, ``$N=2$ Superconformal Models As Topological Field Theories,'' Mod. Phys. Lett. \textbf{A4} (1990) 1653.

% ==== 原文 PDF 第 90 页（印刷页 90）====
\label{op:90}
% [文献 31: Vafa 的拓扑 Landau--Ginzburg 模型；§3.4 中 B 模型在 $r\ll0$ 时把路径积分化归为有限维积分，正是 [31] 所分析的对象。]
\bibitem{ref31} C. Vafa, ``Topological Landau-Ginzburg Models,'' Mod. Phys. Lett. \textbf{A6} (1991) 337.

% [文献 32: Bradlow、Garcia-Prada 等关于稳定对/涡旋模空间的数学文献；§3.4 中把 (3.47) 的前两式（复方程）与第三式（矩映射）的等价解释为“复商 = 辛商”的无穷维类比（Kempf--Ness 定理，§4 中概述）。]
\bibitem{ref32} S. Bradlow, ``Special Metrics And Stability For Holomorphic Bundles With Global Sections,'' J. Diff. Geom. \textbf{33} (1991) 169; S. Bradlow and G. Daskapoulos, ``Moduli Of Stable Pairs For Holomorphic Bundles Over Riemann Surfaces,'' Int. J. Math. \textbf{2} (1991) 477; O. Garcia-Prada, ``Dimensional Reduction Of Stable Bundles, Vortices, And Stable Pairs,'' in preparation.

\bibitem{ref33} V. Guillemin and S. Sternberg, ``Geometric Quantization And Multiplicities Of Group Representations,'' Invent. Math. \textbf{67} (1982) 515.

% [Tutor Note: 文献 [34]、[35] 是几何不变量理论（GIT）的两部标准入门书：Mumford--Fogarty 是 GIT 的原始专著，Newstead 侧重模问题与轨道空间。§4 用初等方式解释“复商 = 辛商”，其严格版本就是 GIT 的稳定点定理（Kempf--Ness），本文 §5.1--§5.4 把加权射影空间/环面簇的构造化归为 GIT 商时反复用到这一语言。]
\bibitem{ref34} D. Mumford and J. Fogarty, \textit{Geometric Invariant Theory}, (Springer, 1982).

\bibitem{ref35} P. Newstead, \textit{Introduction To Moduli Problems And Orbit Spaces} (Tata Institute, 1978).

% [文献 36: Duistermaat--Heckman 积分公式的原始文献；§4 用它说明辛商上由 $\omega|_{D=0}$ 诱导的辛结构类 $[\omega_r]\in H^2(Y//U(1))$ 随 FI 参数 $r$ 线性变化。]
\bibitem{ref36} J. J. Duistermaat and G. J. Heckman, ``On The Variation In The Cohomology In The Symplectic Form Of The Reduced Phase Space,'' Invent. Math. \textbf{69} (1982) 259.

\bibitem{ref37} M. Kreuzer and H. Skarke, ``On The Classification Of Quasi-Homogeneous Functions,'' ``No Mirror Symmetry In Landau-Ginzburg Spectra'' (preprints, 1992).

\bibitem{ref38} M. Aubin, \textit{The Topology Of Torus Action On Symplectic Manifolds} (Birkhauser, 1991).

\bibitem{ref39} D. A. Cox, ``The Homogeneous Coordinate Ring Of A Toric Variety'' (preprint, Amherst College Mathematics Department, 1992).

\bibitem{ref40} V. Batyrev, ``Variations Of The Mixed Hodge Structure Of Affine Hypersurfaces In Algebraic Tori'' (preprint, 1992).

\bibitem{ref41} T. Hubsch, \textit{Calabi-Yau Manifolds: A Bestiary For Physicists} (World Scientific, 1992).

% ==== 原文 PDF 第 91 页（印刷页 91）====
\label{op:91}
\bibitem{ref42} R. Schimmrigk, ``A New Construction Of A Three Generation Calabi-Yau Manifold,'' Phys. Lett. \textbf{193B} (1987) 175.

\bibitem{ref43} D. Gepner, ``Exactly Solvable String Compactification On Manifolds of $SU(N)$ Holonomy,'' Phys. Lett. \textbf{199B} (1987) 380; ``String Theory On Calabi-Yau Manifolds: The Three Generations Case'' PUPT-88/0085 (unpublished); ``Space-Time Supersymmetry In Compactified String Theory And Superconformal Models,'' Nucl. Phys. \textbf{B296} (1988) 757.

% [文献 44: Kollár 关于三维代数簇与 Mori 纲领的综述；§5.5 拓扑改变（flop、收缩）的数学背景。]
\bibitem{ref44} J. Koll\'ar, ``The Structure Of Algebraic Threefolds: An Introduction To Mori's Program,'' Bull. Am. Math. Soc. \textbf{17} (1987) 211.

\bibitem{ref45} R. Hartshorne, \textit{Algebraic Geometry} (Springer-Verlag, 1977).

\bibitem{ref46} M. Dine, N. Seiberg, E. Witten, and X.-G. Wen, ``Non-Perturbative Effects On The String World Sheet I,II,'' Nucl. Phys. \textbf{B278} (1986) 769, \textbf{B289} (1987) 319.

% [文献 47: Candelas–de la Ossa–Green–Parke 的镜像对称周期计算，是本文 $t$ 与复结构模对偶的对照文献。]
\bibitem{ref47} P. Candelas, P. Green, L. Parke, and X. de la Ossa, ``A Pair Of Calabi-Yau Manifolds As An Exactly Soluble Superconformal Field Theory,'' Nucl. Phys. \textbf{B359} (1991) 21.

\bibitem{ref48} P. Aspinwall and D. B. Morrison, ``Topological Field Theory And Rational Curves,'' Oxford preprint, to appear in Commun. Math. Phys.

\bibitem{ref49} C. Hull and E. Witten, ``Supersymmetric Sigma Models And The Heterotic String,'' Phys. Lett. \textbf{160B} (1985) 398.

\bibitem{ref50} M. Dine and N. Seiberg, ``(2,0) Superspace,'' Phys. Lett. \textbf{B180} (1986) 364.

\bibitem{ref51} R. Brooks, J. Gates, and F. Muhammed, Nucl. Phys. \textbf{B268} (1986) 599, ``Extended $D=2$ Supergravity Theories And Their Lower Superspace Realizations,'' Class. Quant. Grav. \textbf{5} (1988) 785.

\bibitem{ref52} M. B. Green, J. H. Schwarz, and E. Witten, \textit{Superstring Theory}, vol. 2 (Cambridge University Press, 1987).

% ==== 原文 PDF 第 92 页（印刷页 92）====
\label{op:92}
\bibitem{ref53} J. Distler and B. Greene, ``Aspects Of (2,0) String Compactifications,'' Nucl. Phys. \textbf{B304} (1988) 1.

\end{thebibliography}
